% from ega4-9-fr.tex
\oldpage[IV-3]{85}

\begin{enumerate}[label=\emph{(\roman*)}]
  \item[(iv)] Pour $i\neq j$, toute composante irréductible de $(Z_i)_s$ est
  distincte de toute composante irréductible de $(Z_j)_s$.
\end{enumerate}
Or, (i) a déjà été vu (9.7.6)~; (ii) est un cas particulier de (9.4.5)~; (iii)
résulte de même de (9.5.3) et (9.4.5). Enfin, si $i\neq j$,
$(Z_i)_\eta\cap(Z_j)_\eta=(Z_i\cap Z_j)_\eta$ est rare dans $(Z_i)_\eta$ ou
dans $(Z_j)_\eta$~; supposons par exemple que $(Z_i\cap Z_j)_\eta$ soit rare
dans $(Z_j)_\eta$~; alors il résulte de (9.5.3) et (9.5.4) que pour $s$ voisin
de $\eta$, $(Z_i\cap Z_j)_s=(Z_i)_s\cap(Z_j)_s$ est rare dans $(Z_j)_s$, ce qui
montre qu'aucune composante irréductible de $(Z_j)_s$ ne peut être contenue
dans une composante irréductible de $(Z_i)_s$, ni \emph{a fortiori} lui être
égale.

Cela étant, il résulte de (ii) et de (3.1.7) que pour $s\in U$, on a
\[
\operatorname{Ass}(\mathcal F_s)\subset
\bigcup_i\operatorname{Ass}((\mathcal F_i)_s),
\]
de (i) que $\operatorname{Ass}((\mathcal F_i)_s)$ est l'ensemble des
$x_{is\alpha}$ ($\alpha\in J_{s,i}$). D'autre part, en vertu de (iii) et du
critère (3.1.3, \emph{c')}), chacun des $x_{is\alpha}$ ($\alpha\in J_{s,i}$,
$i\in I$) appartient à $\operatorname{Ass}(\mathcal F_s)$. Enfin (iv) signifie
que les $x_{is\alpha}$ sont deux à deux distincts.

Reste à prouver la dernière assertion de l'énoncé. Or, il résulte de (i) que
pour $s$ et pour un $i$ donnés, aucun des ensembles $\{x_{is\alpha}\}$ ne peut
être contenu dans un autre pour $\alpha\in J_{s,i}$. D'autre part, si
$(Z_j)_\eta\subset(Z_i)_\eta$, on peut supposer $U$ pris tel que
$(Z_j)_s\subset(Z_i)_s$ pour $s\in U$ (9.5.1), donc chaque $x_{js\beta}$
appartient à l'adhérence d'un $x_{is\alpha}$~; au contraire, si
$(Z_j)_\eta\cap(Z_i)_\eta$ est rare dans $(Z_j)_\eta$, on a vu en prouvant (iv)
que $(Z_j)_s\cap(Z_i)_s$ est rare dans $(Z_j)_s$, donc aucun des $x_{js\beta}$
n'est adhérent à un $x_{is\alpha}$. En particulier, si $(Z_j)_\eta$ est
maximal, ce qui équivaut à dire que $(Z_j)_\eta\cap(Z_i)_\eta$ est rare dans
$(Z_j)_\eta$ pour \emph{tout} $i\neq j$, on en conclut que
$(Z_j)_s\cap(Z_i)_s$ est rare dans $(Z_j)_s$ pour tout $i\neq j$, donc que tout
$x_{js\alpha}$ ($\alpha\in J_{s,j}$) est \emph{maximal}. C.Q.F.D.

\begin{corollary}[9.8.4]
\label{IV.9.8.4-fr}
Les notations et hypothèses étant celles de (9.8.2), il existe un voisinage $U$
de $\eta$ dans $S$ tel que, pour tout $s\in U$, chacun des $(\mathcal F_i)_s$
soit sans cycle premier associé immergé~; en outre, si
$(\mathcal F_{is\alpha})_{\alpha\in J_{s,i}}$ est l'unique décomposition
irredondante réduite de $(\mathcal F_i)_s$, alors, pour tout $s\in U$, la
famille $(\mathcal F_{is\alpha})_{i\in I,\,\alpha\in J_{s,i}}$ est une
décomposition irredondante réduite de $\mathcal F_s$.
\end{corollary}

La première assertion résulte de (9.8.1) et de la définition des
$\mathcal F_i$. D'autre part, on a vu dans (9.8.3) que l'homomorphisme
$\mathcal F_s\to\bigoplus_i(\mathcal F_i)_s$ est injectif et par définition il
en est de même de chacun des homomorphismes
$(\mathcal F_i)_s\to\bigoplus_\alpha\mathcal F_{is\alpha}$, donc
l'homomorphisme $\mathcal F_s\to\bigoplus_{i,\alpha}\mathcal F_{is\alpha}$ est
injectif. Comme on peut supposer (9.4.5) que pour $s\in U$ chacun des
$(\mathcal F_i)_s$ est un quotient de $\mathcal F_s$, les
$\mathcal F_{is\alpha}$ sont des quotients de $\mathcal F_s$, et il reste à
vérifier (3.2.5) que les $x_{is\alpha}$ sont deux à deux distincts et
appartiennent à $\operatorname{Ass}(\mathcal F_s)$, ce qui a été prouvé dans
(9.8.3).

\begin{proposition}[9.8.5]
\label{IV.9.8.5-fr}
Soient $f:X\to S$ un morphisme de présentation finie, $\mathcal F$ un
$\mathcal O_X$-Module quasi-cohérent de présentation finie. Pour tout $s\in S$,
soit $E(s)$ (resp. $E'(s)$) l'ensemble fini (partie de
$\mathbf Z\cup\{-\infty\}$) des dimensions des cycles premiers associés à
$\mathcal F_s$ (resp. des cycles premiers maximaux associés à $\mathcal F_s$,
c'est-à-dire les composantes irréductibles de $\operatorname{Supp}(\mathcal F_s)$).
Alors les fonctions $s\mapsto E(s)$ et $s\mapsto E'(s)$ sont localement
constructibles dans $S$.
\end{proposition}

